\chapter*{Abstract}
\addcontentsline{toc}{chapter}{Abstract}

The exponential growth of the Hilbert space with the number of quantum constituents renders many quantum information tasks intractable for exact classical computation, motivating the development of scalable approximation methods. This thesis presents two physics-informed deep learning frameworks that address such challenges by embedding the underlying physical principles directly into the network architecture and training objective, ensuring that the learned solutions remain consistent with the laws of quantum mechanics.

The first framework concerns the dynamics of closed quantum systems governed by time-dependent Hamiltonians, whose evolution is described by the von Neumann equation. In this setting, the time-evolution operator is expressed through infinite integral expansions. Unlike the time-independent case, where the solution admits a simple exponential form, obtaining an analogous expression requires the introduction of the time-ordering operator. Consequently, when the Hamiltonian does not commute with itself at different times, the mathematical complexity of the problem increases significantly. The proposed method reduces the computational cost of simulating such dynamics. Given the time-evolution operator at a sparse set of sampled time points, a Physics-Informed Neural Network interpolates its evolution over the entire time domain without requiring explicit knowledge of the Hamiltonian. Instead, the model incorporates unitarity as a physical constraint and, for larger systems, predicts an effective Hamiltonian that is propagated using structure-preserving integrators. Across systems ranging from 2 to 8 qubits, the framework achieves mean gate fidelities close to 0.99, accurately reproduces the dynamics even under coarse temporal sampling, and delivers inference speedups of approximately one order of magnitude compared with a geodesic interpolation baseline.

The second framework addresses the Quantum Marginal Problem, which seeks to reconstruct a global quantum state from a prescribed set of reduced density matrices (quantum marginals) while ensuring that all marginals are mutually compatible. The problem becomes increasingly challenging as the system size grows because of the exponential scaling of the Hilbert space and the highly nontrivial compatibility constraints among the marginals. By combining a convolutional denoising autoencoder with the Marginal Imposition Operator, and representing density matrices as two-channel images with a loss function that enforces Hermiticity, positivity, and normalization, the proposed framework reconstructs physically valid global states for systems of 3 to 8 qubits. A hybrid reconstruction scheme recovers states whose marginals match the target marginals with perfect fidelity in many cases, while transfer learning enables the model to generalize across different system sizes with only limited retraining. Once trained, the framework performs inference substantially faster than state-of-the-art semidefinite programming solvers, while successfully reconstructing states in regimes where those methods become computationally impractical.

Taken together, these results demonstrate that incorporating physical principles directly into the learning process enables deep learning to serve not as a replacement for physical modeling, but as its natural extension. By combining data-driven learning with the mathematical structure of quantum mechanics, the proposed frameworks efficiently address problems whose exact classical treatment is computationally prohibitive.